% Control applicability and application in PAIR-AI.
%
% Standalone so it can be compiled and checked on its own; the body from
% \section{...} onward lifts directly into a thesis chapter.
%
% Figures are the two things a reader is most likely to doubt: that applying a
% control is a query rather than code, and that doing so cannot quietly suppress
% the findings it was not aimed at.

\documentclass[11pt,a4paper]{article}

\usepackage[T1]{fontenc}
\usepackage[utf8]{inputenc}
\usepackage{amsmath}
\usepackage{amssymb}
\usepackage{amsthm}
\usepackage{listings}
\usepackage{booktabs}
\usepackage{url}

\theoremstyle{definition}
\newtheorem{definition}{Definition}
\theoremstyle{plain}
\newtheorem{proposition}{Proposition}
\newtheorem{obligation}{Design obligation}

\lstdefinelanguage{SPARQL}{
  morekeywords={CONSTRUCT,WHERE,FILTER,NOT,EXISTS,UNION,OPTIONAL,BIND,VALUES,%
                PREFIX,SELECT,IRI,CONCAT,STR,ENCODE_FOR_URI,AS,a},
  sensitive=true,
  morecomment=[l]{\#},
  morestring=[b]",
}
\lstset{
  language=SPARQL,
  basicstyle=\ttfamily\small,
  keywordstyle=\bfseries,
  commentstyle=\itshape,
  columns=fullflexible,
  breaklines=true,
  frame=single,
  framesep=4pt,
  xleftmargin=6pt,
}

\newcommand{\Lib}{\mathcal{L}}
\newcommand{\Arch}{G}
\newcommand{\Match}{\mathsf{Match}}
\newcommand{\Find}{\mathsf{Find}}
\newcommand{\Prop}{\mathsf{Prop}}
\newcommand{\Assess}{\mathsf{Assess}}
\newcommand{\eval}[2]{[\![#1]\!]_{#2}}

\begin{document}

\section{Applying a control as a graph rewrite}
\label{sec:control-applicability}

An assessment that only reports is a description. What makes a design-time
method actionable is that the analyst can change the design and see the analysis
answer differently. This section defines how PAIR-AI decides that a control is
\emph{applicable} to a particular candidate finding, and how applying it edits
the architecture graph. Both are expressed as registered SPARQL
\texttt{CONSTRUCT} queries, in the same form as motif matching and risk
evaluation, so the knowledge of \emph{where a control belongs} lives beside the
rule that raised the finding rather than in application code.

\subsection{The assessment as a family of construct rules}

\begin{definition}[Assessment state]
Let $\Lib$ be the curated library: motifs, risk patterns, controls, and the
declarations that register their implementations. Let $\Arch$ be a submitted
architecture graph. Every executable rule is a \texttt{CONSTRUCT} query $q$,
and $\eval{q}{H}$ denotes the set of triples $q$ constructs over a graph $H$.
\end{definition}

Rules are partitioned by a declared output type,
$\mathsf{out}(q) \in \{\Prop, \Match, \Find, \mathsf{Mit}\}$, corresponding to
\texttt{pair:DataCategoryPropagation}, \texttt{pair:MotifMatch},
\texttt{pair:RiskFinding} and \texttt{pair:MitigationApplication}. Writing $Q_X$
for the rules of type $X$, the pipeline is

\begin{align}
  \Prop(\Arch)  &= \operatorname{lfp}\Big(H \mapsto H \cup \textstyle\bigcup_{q \in Q_\Prop} \eval{q}{H}\Big)(\Lib \cup \Arch), \label{eq:prop}\\
  \Match(\Arch) &= \textstyle\bigcup_{q \in Q_\Match} \eval{q}{\Prop(\Arch)}, \label{eq:match}\\
  \Find(\Arch)  &= \textstyle\bigcup_{q \in Q_\Find} \eval{q}{\Prop(\Arch) \cup \Match(\Arch)}, \label{eq:find}\\
  \Assess(\Arch) &= \Prop(\Arch) \cup \Match(\Arch) \cup \Find(\Arch). \label{eq:assess}
\end{align}

Equation~\eqref{eq:prop} is a least fixed point because a propagated data
category can satisfy the precondition of another propagation rule; the
implementation iterates to stability under a fixed bound. Rules of type
$\mathsf{Mit}$ appear nowhere in \eqref{eq:prop}--\eqref{eq:assess}. That
omission is deliberate and is the subject of
Proposition~\ref{prop:no-self-mitigation}.

\subsection{Monotonicity, and why a control can clear a finding}

Motif matching is monotone: a motif asserts the presence of structure, never its
absence, so extending the graph cannot destroy a match.

\begin{proposition}[Monotone matching]
\label{prop:monotone-match}
If $\Arch \subseteq \Arch'$ then $\Match(\Arch) \subseteq \Match(\Arch')$.
\end{proposition}

\noindent Risk evaluation is different. An applicability condition may test for
the \emph{absence} of a represented control, which in SPARQL is
\texttt{FILTER NOT EXISTS} and is therefore closed-world over the submitted
graph. Consequently $\Find$ is \emph{not} monotone in $\Arch$: adding triples
can remove a finding.

That non-monotonicity is not a defect to be tolerated; it is precisely the
property a mitigation exploits. If risk evaluation were monotone, no addition to
an architecture could ever discharge a finding, and the only way to clear one
would be to assert that it had been handled --- which is a claim about the
analyst, not about the design.

\begin{obligation}[Closed-world reading]
\label{obl:cwa}
Because the negation is closed-world over $\Arch$, the disappearance of a
finding means \emph{the amended graph no longer represents the gap}. It does not
mean the risk has been eliminated. Every finding is a candidate throughout, and
the wording of the tool must not drift from that reading.
\end{obligation}

\subsection{Applicability}

Each finding records the rule that produced it, via
\texttt{pair:\allowbreak generatedByRiskPattern}; write $\rho(f)$ for that risk
pattern. Each mitigation rule $q$ declares the control it realises,
$\mathsf{ctl}(q)$ via \texttt{pair:implementsControl}, and the risk pattern it
answers, $\mathsf{rp}(q)$ via \texttt{pair:mitigatesRiskPattern}.

\begin{definition}[Applicable control]
\label{def:applicable}
A control $c$ is \emph{applicable} to a finding $f$ iff
\[
  \exists\, q \in Q_{\mathsf{Mit}} \; : \;
  \mathsf{ctl}(q) = c \;\wedge\; \mathsf{rp}(q) = \rho(f).
\]
\end{definition}

Applicability is thus a property of the \emph{pair} $(c, \rho(f))$, not of the
control alone. The distinction is load-bearing. A single control is suggested by
several risk patterns --- output validation answers improper output handling and
sensitive information disclosure --- while a rewrite is written against one
vulnerable shape. Keyed on the control alone, every finding suggesting that
control would offer the same rewrite; the rewrite would then evaluate against a
shape it was not written for, construct nothing, and report that the control was
already in place while the risk remained. Definition~\ref{def:applicable} is what
prevents that.

A control that is not applicable is still reported as a suggestion. Many
controls are non-technical --- provenance, vetting, governance --- and have no
runtime shape at all; for those, no rewrite exists and none should be invented.

\subsection{Application}

\begin{definition}[Application]
\label{def:apply}
Let $c$ be applicable to $f$ through rule $q$, and let
$H = \Assess(\Arch)$ be the assessed graph, which is where $f$ exists. The
\emph{insertion} is
\[
  \Delta(f, c) \;=\; \eval{q}{H}\big[\, \mathtt{?finding} \mapsto f \,\big],
\]
the substitution binding the query variable \texttt{?finding} to $f$ before
evaluation. The amended architecture is $\Arch^{+} = \Arch \cup \Delta(f,c)$.
\end{definition}

Two aspects of Definition~\ref{def:apply} deserve comment.

\paragraph{The insertion is scoped to one finding.} Binding \texttt{?finding}
externally, rather than letting the rule range over all findings, is what makes
this an analyst's action on a chosen problem instead of a blanket edit.

\paragraph{The insertion lands in $\Arch$, not in $H$.} Only the architecture is
amended; the library is never written to, and the derived facts of the previous
run are discarded. The analyst then re-runs \eqref{eq:assess} over $\Arch^{+}$,
and the same rules speak again. The tool does not declare the finding resolved:
it edits the design and recomputes.

\subsection{The shape of a mitigation rule}

A mitigation rule restates the vulnerable configuration its
$\mathsf{rp}(q)$ detected, and constructs the element that interrupts it, bound
to the elements the finding already cites. Nothing is inferred about which
evidence element plays which part, because the conditions that raised the
finding also identify where the control belongs. The rule answering prompt
injection is representative:

\begin{lstlisting}
CONSTRUCT {
    ?screeningStep a beam:Process ;
        pair:playsRole pair:InputGuardrailStep ;
        beam:use    ?untrustedContent ;
        beam:inform ?generationStep .
    ?system beam:hasProcess ?screeningStep .
}
WHERE {
    ?finding a pair:RiskFinding ;
        pair:generatedByRiskPattern pat:PromptInjectionRiskPattern ;
        pair:hasEvidence ?untrustedContent , ?generationStep .

    # the vulnerable shape, restated
    ?untrustedContent pair:containsDataCategory/pair:subDataCategoryOf*
                      pair:UntrustedContent .
    ?generationStep pair:playsRole/pair:subRoleOf* pair:GenerationStep .
    { ?generationStep beam:use ?untrustedContent }
    UNION
    { ?intermediateStep beam:use ?untrustedContent ;
                        beam:produce ?derivedContext .
      ?generationStep beam:use ?derivedContext }

    # and only where no screen is represented yet
    FILTER NOT EXISTS {
        ?existingScreen pair:playsRole/pair:subRoleOf* pair:InputGuardrailStep .
        { ?existingScreen beam:use ?untrustedContent ;
                          beam:inform ?generationStep }
        UNION
        { ?existingScreen beam:use ?untrustedContent ;
                          beam:produce ?screened .
          ?generationStep beam:use ?screened }
    }

    BIND(IRI(CONCAT(
        "http://w3id.org/airiskkg/generated/mitigation/input-screening/",
        ENCODE_FOR_URI(STR(?untrustedContent)), "/",
        ENCODE_FOR_URI(STR(?generationStep))
    )) AS ?screeningStep)
}
\end{lstlisting}

Applied to a direct-prompting architecture, this constructs six triples. Every
subject and object other than the newly minted step is an element already
present in $\Arch$, which is what distinguishes an edit from an unconnected
addition: one new node, wired to the input it screens and the generation it
informs.

\subsection{Properties}

\begin{proposition}[No self-mitigation]
\label{prop:no-self-mitigation}
For any $\Arch$, $\Assess(\Arch)$ is independent of $Q_{\mathsf{Mit}}$.
\end{proposition}

\begin{proof}
Immediate from \eqref{eq:prop}--\eqref{eq:assess}: the pipeline selects rules by
declared output type and requests only $\Prop$, $\Match$ and $\Find$. No rule of
type $\mathsf{Mit}$ is evaluated.
\end{proof}

\noindent This separation is the safety catch of the design. Were mitigation
rules evaluated inside the pipeline, each would construct the very structure its
risk pattern tests for the absence of, every finding would discharge itself, and
the assessment would report nothing. The property is enforced by a regression
test rather than left to convention.

\begin{proposition}[Idempotence]
\label{prop:idempotent}
$\Delta\big(f, c\big)$ recomputed over $\Assess(\Arch^{+})$ is empty.
\end{proposition}

\begin{proof}[Argument]
Two independent reasons. First, the guard: the rule's
\texttt{FILTER NOT EXISTS} tests for the same structure the rule constructs, so
after insertion the guard fails and no solution survives. Second, minting: the
inserted IRI is a deterministic function of the elements it is bound to, so even
were the guard to admit a solution, the constructed triples would be identical
to those already present and their union with $\Arch^{+}$ would add nothing.
\end{proof}

\begin{obligation}[Discharge]
\label{obl:discharge}
For each rule $q$, the structure $q$ constructs must satisfy the
absence-condition of $\mathsf{rp}(q)$, so that $f \notin \Find(\Arch^{+})$.
\end{obligation}

\noindent Obligation~\ref{obl:discharge} is a property of each rule against its
risk pattern, not a theorem about the framework, and is discharged by test per
rule. It is stated separately because the converse failure is the one that
matters: a rule may construct something plausible that does not satisfy the
condition, in which case the analyst edits the design and the finding persists.

\begin{obligation}[No collateral discharge]
\label{obl:collateral}
Applying a control must not remove findings other than by mitigating them.
\end{obligation}

\noindent This obligation is easy to violate. An interposing rewrite --- one
that rewires a flow so the control sits \emph{on} the path rather than reading
from it --- was implemented and withdrawn during development: inserting an
output screen removed prompt-injection and system-prompt-leakage findings,
because several risk patterns require a generation step to produce a user-facing
output \emph{directly}, and interposition breaks that pattern rather than the
risk. The untrusted content still reached the prompt; the analysis simply
stopped reporting it. For a method whose outputs are candidate risks, a false
negative introduced by a mitigation is worse than a coarse-grained one, and the
prerequisite --- path-tolerant risk conditions --- is a separate change.

\subsection{Coverage}

\begin{table}[ht]
\centering
\begin{tabular}{@{}llr@{}}
\toprule
Rule type & Declared output type & Count \\
\midrule
Propagation        & \texttt{pair:DataCategoryPropagation} & 6 \\
Motif matching     & \texttt{pair:MotifMatch}              & 31 \\
Risk evaluation    & \texttt{pair:RiskFinding}             & 15 \\
Mitigation rewrite & \texttt{pair:MitigationApplication}   & 9 \\
\bottomrule
\end{tabular}
\caption{Registered rules by declared output type. Only the first three are
evaluated during an assessment (Proposition~\ref{prop:no-self-mitigation}).}
\label{tab:rules}
\end{table}

Twelve of the fifteen risk patterns admit a structural discharge; the remaining
three --- data and model poisoning, supply-chain compromise, and vector and
embedding weakness --- rest on provenance and vetting, whose controls are
non-technical and have no runtime shape. For those, no rewrite can exist, and
the appropriate mechanism is analyst triage recorded against the finding rather
than a fabricated structural fact in the graph. Conflating the two was a
concrete defect in an earlier version: fifteen risk rules carried an escape
clause satisfied by an annotation that no component of the system ever wrote, so
findings appeared falsifiable while the only thing that could discharge them was
unreachable.

\end{document}
